%% The line mask: the list, the frame, what it costs, and the three things a
%% reader has to be able to check -- the dates, the recurrence test applied to
%% the labelled crossings, and the image sideband.
%% Owner: r12-mask.  Carries \label{app:linemask} (cited by 04_method.tex and
%% app_C_spatialnorm.tex) and \label{tab:masklines}.
\section{The line mask}
\label{app:linemask}

A crossing that falls on a known molecular transition has an astrophysical
explanation to hand, and the mask exists to say which crossings those are. It
is a label and not a verdict: every crossing goes on to the spatial and
recurrence tests whether or not it carries the label, and the labelled
crossings are tested below.

\emph{The list, and when it was made.} The mask holds every transition of the
\FrNSpec{} species these spectral setups cover --- the three CO isotopologues,
HCN, HCO$^{+}$, CS, CN, SiO, H$_2$CO, SO, [C\,\textsc{i}] and hydrogen
recombination --- that falls inside one of the \McNIsland{} searched frequency
intervals with an upper-state energy below \McEuMax\,K. That rule is the whole
of it, and it admits \FrNTrans{} transitions
(Table~\ref{tab:masklines}). No strength cut is applied across species,
because the Einstein coefficients of the CO ladder are
$10^{-7}$--$10^{-6}$\,s$^{-1}$ and any threshold admitting the strong lines of
HCN would discard the brightest lines in these fields; within a species the
coefficient keeps a rotational transition with its main hyperfine components
and drops the weak satellites. The rarer isotopologues, $^{34}$SO and
H$^{13}$CN and the rest, are not among the named species and would add
\McNRare{} transitions. The list is a single Splatalogue query of CDMS and JPL
\citep{Pickett1998,CDMS2005} made on \MkListDate{} and frozen with the data
release, so it is reproducible as a file rather than as a description --- and
it therefore post-dates the \MkCritDate{} fixing of the disposition criteria,
as does every species in it. It is a uniform rule applied to all \NCross{}
crossings, not a record of what was recognised when each was first examined.

\emph{The frame is the star's own, and the sign is not the usual one.} A
crossing's frequency is carried from the topocentric value to the barycentre
using its own block's pointing and mid-time, and then to the stellar rest
frame using the star's systemic velocity. The barycentric term runs from
\FrBaryLo{} to \FrBaryHi\,km\,s$^{-1}$ across the survey, so a tolerance
applied in the observed frame would spend \FrBarySwing\,km\,s$^{-1}$ of its
budget on the Earth's motion. The \FrEpochNEpoch{} \FrEpochStar{} epochs show
the transform working in the only way that can be checked:
\FrEpochTopoKms\,km\,s$^{-1}$ apart as observed, they agree to
\FrEpochStelKms\,km\,s$^{-1}$ in the star's frame. Throughout,
$\Delta v_\star = c\,(\nu_\star-\nu_0)/\nu_0$, so a positive offset means the
crossing lies \emph{above} the transition in frequency --- the opposite sign
to the velocity convention of a spectral line.

\emph{The half-width, and what it costs.} A one-channel test is too tight for
circumstellar gas, whose emission is shifted by many channels: the tolerance
follows from Keplerian motion where these molecules survive,
$\lesssim$13\,km\,s$^{-1}$ beyond 10\,au, and is
$\pm\MaskHalfKms$\,km\,s$^{-1}$. At that width the mask removes
\FrMaskLostGHz\,GHz of the \FrMaskUnionGHz\,GHz unique-frequency union, which
is \FrMaskPct{} per cent of it; the loss is concentrated in the line-tuned
fine-channel windows, and the cost band by band is in the data release.
Between $\pm\FrLadderLo$ and $\pm\FrLadderHi$\,km\,s$^{-1}$ the mask labels
\FrAttrWTwenty{} to \FrAttrWHundred{} of the \NCross{} crossings, and the
\FrNRankPass{} unattributed crossing that outranks all \NCtrl{} control
positions in its own window is the same one at every width.

\begin{table}
\centering
\footnotesize
\caption{The line mask, by species: the number of transitions $N$ it holds,
the frequency range they span inside the searched intervals and their
upper-state energies. Hydrogen recombination lines carry no tabulated
$E_{\rm u}$. The transition-by-transition list is in the data release.}
\label{tab:masklines}
%% GENERATED by maskrecur_v413.py -- do not hand-edit.
\begin{tabular}{@{}lrr@{\,}c@{\,}rr@{}}
\hline
Species & $N$ & \multicolumn{3}{c}{$\nu_0$ (GHz)} & $E_{\rm u}$ (K) \\
\hline
H & 5 & 92.03 & -- & 231.90 & -- \\
SO & 15 & 100.03 & -- & 682.68 & 16--143 \\
H$_2$CO & 13 & 101.33 & -- & 491.97 & 10--106 \\
C$^{18}$O & 3 & 109.78 & -- & 329.33 & 5--32 \\
CN & 50 & 113.12 & -- & 680.30 & 5--114 \\
CO & 3 & 115.27 & -- & 345.80 & 6--33 \\
CS & 3 & 195.95 & -- & 342.88 & 24--66 \\
SiO & 2 & 217.10 & -- & 347.33 & 31--75 \\
$^{13}$CO & 8 & 220.40 & -- & 330.59 & 16--32 \\
HCN & 7 & 354.50 & -- & 354.51 & 43--43 \\
HCO$^{+}$ & 1 & 356.73 & -- & 356.73 & 43 \\
{[C\,\textsc{i}]} & 1 & 492.16 & -- & 492.16 & 24 \\
\hline
\end{tabular}%

\end{table}

\emph{Sulphur monoxide decides \MkNSODecideWord{} coincidences, and one rule is
applied to all of them.} \MkSOStars{} lie \MkSODvList\,km\,s$^{-1}$ from sulphur monoxide
transitions of upper-state energy \MkSOEuLo--\MkSOEuHi\,K, each inside a
Band~7 window that also covers CO($3\to2$), and none has a nearer transition
of any other masked species. There is reason for caution about each:
sulphur monoxide has not been reported in a debris disc, where the
second-generation gas is CO, C and O, and toward \FrCpStar{} the CO($3\to2$)
line is covered in the same window and absent to \FrCpCoLimMJy\,mJy in a
\FrCpCoChanKms\,km\,s$^{-1}$ channel. That caution applies equally to each of
them, and so does the date, since the species list post-dates every
classification; neither separates one coincidence from the others. Applying
the rule attributes all of them and gives \MkAttrSO{} attributed crossings
against \MkUnattrSO{} unattributed; withholding sulphur monoxide from all of
them instead gives \MkAttrNoSO{} and \MkUnattrNoSO, and returns a second
crossing to the \MkRankNoSO{} that outrank every control in their own
windows. The
recurrence test says the same thing about each of them either way, so the choice
moves counts and no conclusion.

\emph{The labelled crossings are tested for recurrence too.} Of the
\MkAttrSO{} attributed crossings, \MkAttrTested{} have a later execution block
covering the predicted cell. \MkAttrRecur{} are present again: toward
\MkAttrRecurStars{} a later block carries its own threshold crossing at the
same stellar-frame frequency, at $T_\star=\MkAttrRecurTLo$ to
\MkAttrRecurTHi, which is the test working in the positive direction on
emission that is certainly real. The other \MkAttrExcl{} are not present
again, and in each the repeat coverage is deep enough that a carrier at the
discovery flux would have exceeded the trigger there and does not; among them
are the sulphur monoxide coincidences, which the test limits to
\MkAttrExclLo--\MkAttrExclHi$\sigma$ whichever side of the mask they are
placed. \MkAttrUntest{} cannot be tested and are
named rather than counted: \MkAttrUntestNames, each observed once at the
frequency concerned.

\emph{Image-sideband leakage does not account for any crossing.} The mask
holds signal-sideband frequencies, so a bright line in the image sideband
could in principle be mistaken for a carrier. \MkImgN{} unattributed
crossings invite the question, at \MkImgFreqs\,GHz toward \MkImgStar, because
each shares a block with a CO($2\to1$) crossing at \MkImgPartner\,GHz in the
opposite sideband. A mirror about the first local oscillator would place that
line on the crossing only for an oscillator at
\MkImgLOReqLo--\MkImgLOReqHi\,GHz, and the block's own spectral windows
require \MkImgLOLo--\MkImgLOHi\,GHz: at the setting the mirror needs, the
lowest window of the block would sit at an intermediate frequency of
\MkImgIFReq\,GHz, beyond the \MkImgIFMax\,GHz the receiver provides
\citep{ALMAHandbook}. Their image frequencies fall at least
\MkImgDvMin\,km\,s$^{-1}$ from every masked transition.

\emph{A wider list is not a better one.} Every Splatalogue transition over the
searched intervals under the same energy cut and with an Einstein coefficient
above $10^{-5}$\,s$^{-1}$ --- \MaskDbTrans{} transitions of \MaskDbCarriers{}
species --- masks \MaskDbPctUnion{} per cent of the union and moves
\FrFullNFlip{} of the \NCross{} crossings across the line,
\FrFullNFlipRankWord{} of them a crossing that outranks every control, at
$\FrFullFlipRankDv$\,km\,s$^{-1}$ from a \FrFullFlipRankChem{} transition. At
these frequencies most frequencies lie within $\pm\MaskHalfKms$\,km\,s$^{-1}$
of something in such a list, which is why the species are restricted to what
these spectral setups cover, and why the outcome is a label. ``Coincident
with a known molecular line'' is a statement about a catalogue rather than a
property of a frequency.
